\documentclass[11pt]{article}

% ===== 基本宏包：尽量保持简单 =====
\usepackage[margin=1in]{geometry}
\usepackage{graphicx}
\usepackage{float}
\usepackage{chngcntr}
\usepackage{booktabs}
\usepackage{multirow}
\usepackage{longtable}
\usepackage{array}
\usepackage{calc}
\usepackage{amsmath, amssymb}
\setlength{\LTcapwidth}{\textwidth}
\usepackage{xurl}
\usepackage{hyperref}
\usepackage[numbers,sort&compress]{natbib}
\usepackage{setspace}
\usepackage{fontspec}
\usepackage{xeCJK}
\providecommand{\tightlist}{\setlength{\itemsep}{0pt}\setlength{\parskip}{0pt}}

% ===== 图表编号风格：全文连续编号，使用 Figure / Table =====
\renewcommand{\figurename}{Figure}
\renewcommand{\tablename}{Table}

% ===== 字体与换行策略 =====
\defaultfontfeatures{Ligatures=TeX,Scale=MatchLowercase}
\IfFontExistsTF{Noto Serif}{\setmainfont{Noto Serif}}{\IfFontExistsTF{DejaVu Serif}{\setmainfont{DejaVu Serif}}{\setmainfont{Latin Modern Roman}}}
\IfFontExistsTF{Noto Sans}{\setsansfont{Noto Sans}}{\IfFontExistsTF{DejaVu Sans}{\setsansfont{DejaVu Sans}}{\setsansfont{Latin Modern Sans}}}
\IfFontExistsTF{Noto Sans Mono}{\setmonofont{Noto Sans Mono}}{\IfFontExistsTF{DejaVu Sans Mono}{\setmonofont{DejaVu Sans Mono}}{\setmonofont{Latin Modern Mono}}}
\IfFontExistsTF{UnBatang}{\setCJKmainfont[AutoFakeBold=3,AutoFakeSlant=.2]{UnBatang}}{\IfFontExistsTF{Baekmuk Batang}{\setCJKmainfont[AutoFakeBold=3,AutoFakeSlant=.2]{Baekmuk Batang}}{\IfFontExistsTF{Noto Serif CJK KR}{\setCJKmainfont[AutoFakeBold=3,AutoFakeSlant=.2]{Noto Serif CJK KR}}{\IfFontExistsTF{Noto Serif CJK SC}{\setCJKmainfont[AutoFakeBold=3,AutoFakeSlant=.2]{Noto Serif CJK SC}}{\setCJKmainfont[AutoFakeBold=3,AutoFakeSlant=.2]{WenQuanYi Zen Hei}}}}}
\IfFontExistsTF{UnDotum}{\setCJKsansfont[AutoFakeBold=3,AutoFakeSlant=.2]{UnDotum}}{\IfFontExistsTF{Baekmuk Dotum}{\setCJKsansfont[AutoFakeBold=3,AutoFakeSlant=.2]{Baekmuk Dotum}}{\IfFontExistsTF{Noto Sans CJK KR}{\setCJKsansfont[AutoFakeBold=3,AutoFakeSlant=.2]{Noto Sans CJK KR}}{\IfFontExistsTF{Noto Sans CJK SC}{\setCJKsansfont[AutoFakeBold=3,AutoFakeSlant=.2]{Noto Sans CJK SC}}{\setCJKsansfont[AutoFakeBold=3,AutoFakeSlant=.2]{WenQuanYi Zen Hei}}}}}
\IfFontExistsTF{UnTaza}{\setCJKmonofont[AutoFakeBold=3]{UnTaza}}{\IfFontExistsTF{Noto Sans Mono CJK KR}{\setCJKmonofont[AutoFakeBold=3]{Noto Sans Mono CJK KR}}{\IfFontExistsTF{Noto Sans Mono CJK SC}{\setCJKmonofont[AutoFakeBold=3]{Noto Sans Mono CJK SC}}{\setCJKmonofont[AutoFakeBold=3]{WenQuanYi Zen Hei Mono}}}}
\setlength{\emergencystretch}{6em}
\sloppy

% 可选：让行距略宽一些，便于审稿阅读
\onehalfspacing

\title{\textbf{Decathlon VOC Analyzer: 상품 이미지·텍스트 증거와 사용자 리뷰 정렬을 위한 증거 기반 멀티모달 VOC 분석 시스템}}

\author{{\small
\begin{tabular}{@{}ccc@{}}
\multicolumn{3}{c}{\parbox[t]{0.6600\textwidth}{\centering \textbf{Severin Ye}\textsuperscript{1} \\ \textit{School of Computer Science and Engineering} \\ \textit{Kyungpook National University} \\ Daegu, Republic of Korea \\ {\ttfamily\small 6severin9\allowbreak@gmail.\allowbreak{}com}}} \\
\multicolumn{3}{c}{} \\[0.9em]
\parbox[t]{0.3067\textwidth}{\centering \textbf{Dokeun Lee}\textsuperscript{2} \\ \textit{Department of Physics} \\ \textit{Kyungpook National University} \\ Daegu, Republic of Korea \\ {\ttfamily\small nsa08008\allowbreak@naver.\allowbreak{}com}} & \parbox[t]{0.3067\textwidth}{\centering \textbf{Wushuang Liu}\textsuperscript{2} \\ \textit{Department of Child Studies} \\ \textit{Kyungpook National University} \\ Daegu, Republic of Korea \\ {\ttfamily\small 824309203q\allowbreak@gmail.\allowbreak{}com}} & \parbox[t]{0.3067\textwidth}{\centering \textbf{Seowan Jung}\textsuperscript{2} \\ \textit{School of Computer Science and Engineering} \\ \textit{Kyungpook National University} \\ Daegu, Republic of Korea \\ {\ttfamily\small swan041014\allowbreak@gmail.\allowbreak{}com}} \\[0.9em]
\multicolumn{3}{c}{\makebox[0.9200\textwidth][c]{\parbox[t]{0.3067\textwidth}{\centering \textbf{HyunJun Jung}\textsuperscript{2} \\ \textit{School of Computer Science and Engineering} \\ \textit{Kyungpook National University} \\ Daegu, Republic of Korea \\ {\ttfamily\small beanbeansummon\allowbreak@naver.\allowbreak{}com}}\hspace{2\tabcolsep}\parbox[t]{0.3067\textwidth}{\centering \textbf{Hye Jeon}\textsuperscript{2} \\ \textit{School of Computer Science and Engineering} \\ \textit{Kyungpook National University} \\ Daegu, Republic of Korea \\ {\ttfamily\small prajna0426\allowbreak@naver.\allowbreak{}com}}}} \\
\end{tabular}
}}

\date{} % Nature 系列投稿一般不需要显示日期

\begin{document}
\maketitle
\nocite{lewis2020retrievalAugmentedGeneration,gao2023rarr,radford2021clip,liu2024structuredOutputBenchmark,geng2023grammarConstrainedDecoding,yao2023react,nvidia2024multimodalRagIntro,nvidia2024aiReadyKnowledgeSystems,faysse2024colpali,pontiki2014semevalAspect,qdrant2024pdfRetrievalAtScale,qdrant2024multiVectorCourse,manning2008introductionIr}

% =========================
% 초록
% =========================
\section{초록}
전자상거래 VOC(Voice of Customer) 분석은 리뷰에 담긴 주관적 사용자 경험과 상품 페이지의 객관적 증거를 함께 이해해야 한다. 본 논문은 증거 기반 멀티모달 VOC 분석 시스템인 Decathlon VOC Analyzer를 제안한다. 이 방법은 상품 정보를 증거 패키지로 표준화하고, 리뷰에서 측면 수준의 신호를 추출하며, 각 측면을 명확한 검색 경로를 가진 질문으로 계획한다. 이후 두 단계의 멀티모달 후보 검색과 재순위를 수행하고, 최종적으로 claim 수준의 증거 귀속을 포함한 구조화 보고서를 생성한다. 시스템은 재현성과 오류 분석을 지원하기 위해 핵심 중간 표현을 보존한다. 본 논문의 기여는 ``리뷰 측면 모델링 - 질문 계획 - 멀티모달 검색 - 증거 제약 생성'' 프레임워크, 구조화 중간 표현에 기반한 재현 가능한 시스템 구현, 그리고 166개 자동화 테스트를 통한 시스템 안정성 검증이다. 본 논문은 시스템 방법 논문으로 자리매김하며, 대규모 수작업 주석 benchmark를 완료했다고 주장하지 않는다.

\subsection{키워드}

멀티모달 VOC 분석; 증거 제약 생성; 이미지-텍스트 검색; 상품 리뷰 마이닝; claim 귀속

% =========================
% I. 서론
% =========================
\section{I. 서론}
사용자 리뷰 분석은 지능형 전자상거래 운영의 기초 과제이다. 기존 방법은 주로 감성 분류, 주제 군집화, 측면 기반 감성 분석에 초점을 두어 사용자가 상품 속성에 대해 보이는 긍정적 또는 부정적 태도를 요약할 수 있다. 그러나 이러한 리뷰 의견이 상품 자체의 증거로 설명, 지지 또는 반박될 수 있는지는 상대적으로 덜 다룬다. 상품 운영과 제품 개선의 관점에서 ``사용자가 용량이 작다고 불평한다''는 사실만으로는 충분하지 않다. 분석자는 상품 페이지가 치수를 명확히 설명하는지, 이미지가 공간 구조를 보여 주는지, 해당 문제가 상품 결함인지 사용자 기대의 차이인지도 알아야 한다.

대규모 언어 모델은 상품 리뷰 요약에 새로운 가능성을 제공하지만, 단일 단계 요약에는 분명한 위험이 있다. 시스템이 상품 문구, 이미지 설명, 리뷰 전체를 그대로 모델에 입력하면 모델은 유창하지만 검토하기 어려운 결론을 생성할 수 있다. 사용자는 특정 제안이 어떤 리뷰, 어떤 상품 텍스트, 어떤 이미지에서 비롯되었는지 판단하기 어렵고, 오류가 리뷰 이해, 검색, 증거 해석, 최종 생성 중 어느 단계에서 발생했는지도 찾기 어렵다. 따라서 상품 VOC 분석은 블랙박스 요약 과제로 축소되어서는 안 되며, 리뷰 요구에서 상품 증거로, 다시 구조화 보고서로 이어지는 추적 가능한 흐름으로 모델링되어야 한다.

본 논문의 연구 문제는 전자상거래 상품 장면에서 사용자 리뷰의 주관적 경험을 상품 문구와 상품 이미지의 검증 가능한 증거에 어떻게 명시적으로 정렬하고, 검토 가능하며 재현 가능하고 평가 가능한 VOC 인사이트를 생성할 것인가이다. 이를 위해 본 논문은 Decathlon VOC Analyzer를 제안한다. 시스템은 상품 중심의 표준화 증거 패키지를 구성하고, 리뷰를 측면 수준 객체로 추출하며, 질문 계획 계층을 통해 각 측면을 텍스트, 이미지 또는 이미지-텍스트 결합 검색 과제로 변환한다. 마지막으로 멀티모달 후보 검색, 재순위, 증거 집계, claim 수준 귀속을 통해 구조화 보고서를 생성한다.

% =========================
% II. 방법
% =========================
\section{II. 방법}
본 논문은 상품 VOC 분석을 증거 기반 방법 절차로 구성한다. 핵심 목표는 리뷰 요약을 직접 생성하는 것이 아니라, 리뷰 속 주관적 피드백을 검색 가능하고 검증 가능하며 귀속 가능한 이미지-텍스트 증거 체인으로 변환하고, 통일된 구조 아래에서 질문 계획, 증거 검색, 실험 설명, 그림 제시를 수행하는 것이다.

\subsection{전체 구조}

Decathlon VOC Analyzer는 상품 VOC 분석을 증거 기반 다단계 프로세스로 형식화한다. 시스템은 먼저 상품 페이지의 텍스트, 이미지, 리뷰를 통일된 증거 공간으로 조직한다. 이후 리뷰 측면에서 측면 수준 사용자 피드백을 추출한다. 다음으로 질문 계획은 주관적 피드백을 실행 가능한 증거 질의로 변환한다. 마지막으로 상품 텍스트와 상품 이미지 사이에서 멀티모달 검색, 재순위, 증거 제약 생성을 수행한다 \citep{lewis2020retrievalAugmentedGeneration,gao2023rarr,radford2021clip,nvidia2024multimodalRagIntro,nvidia2024aiReadyKnowledgeSystems,faysse2024colpali}. 단일 단계 대규모 모델 요약과 달리, 이 프레임워크는 ``사용자가 무엇을 표현했는가'', ``상품 증거가 이를 지지할 수 있는가'', ``최종 보고서가 어떻게 귀속되는가''를 관찰 가능한 중간 단계로 분해한다.

시스템은 네 개의 논리 계층으로 구성된다. 상품 증거 계층은 상품 텍스트, 이미지, 국소 시각 영역의 구조화 표현을 구축한다. VOC 요구 계층은 리뷰를 측면, 의견, 감성, 사용 장면으로 변환한다. 증거 정렬 계층은 측면을 명확한 모달리티 경로를 가진 질문으로 다시 쓰고, 상품 증거 공간에서 후보를 검색하고 재순위화한다. 보고서 귀속 계층은 구조화 인사이트를 생성하고 보고서 claim을 리뷰 증거와 상품 증거에 매핑한다.

\subsection{수식}

리뷰 측면 \(a\)와 후보 증거 \(e\)의 순위를 통일하기 위해, 본 논문은 텍스트 관련성, 이미지 관련성, 교차 모달 일관성을 하나의 증거 점수로 결합한다.

\[
S(a, e) = \lambda_t R_t(a, e) + \lambda_v R_v(a, e) + \lambda_c C(a, e)
\]

여기서 \(R_t(a, e)\)는 텍스트 경로의 관련성 점수, \(R_v(a, e)\)는 이미지 경로의 관련성 점수, \(C(a, e)\)는 리뷰 측면과 상품 증거 사이의 교차 모달 일관성 제약을 의미한다. \(\lambda_t\), \(\lambda_v\), \(\lambda_c\)는 조정 가능한 가중치이다. 이 식은 최종 생성 분석을 대체하기 위한 것이 아니라, 검색과 재순위 단계에서 통일된 증거 선택 기준을 제공하기 위한 것이다.

\subsection{실험 설명}

상품 증거 표현 단계에서 시스템은 상품 제목, 설명, 카테고리, 규격, 이미지, 사용자 리뷰를 안정적인 식별자를 가진 증거 객체로 표현한다. 텍스트 증거는 출처 필드를 보존하고, 이미지 증거는 전체 이미지와 규칙 기반 국소 영역을 함께 보존하며, 리뷰 증거는 평점, 언어 힌트, 원문 조각을 보존한다. 이를 통해 이후 검색, 귀속, 평가를 위한 통일된 객체 경계를 마련한다.

리뷰 모델링 단계에서 시스템은 자연어 피드백을 측면 수준 VOC 단위로 변환한다 \citep{pontiki2014semevalAspect}. 각 측면 단위는 평가 대상 속성, 감성 경향, 의견 설명, 원문 증거 조각, 사용 장면, 신뢰도를 포함한다. 측면 수준 표현은 통계 집계뿐 아니라 특정 속성을 중심으로 한 증거 질의 생성도 지원한다. 제한된 리뷰 예산이 높은 별점의 짧은 리뷰에 점유되는 것을 피하기 위해, 시스템은 리뷰 추출 전에 별점 인지 샘플링을 도입하고 구조화 LLM 경로와 휴리스틱 경로를 모두 지원하여 모델 능력과 프로세스 완전성을 구분한다.

질문 계획 단계에서 시스템은 원시 리뷰를 직접 사용해 상품 증거를 검색하지 않는다. 대신 각 측면을 명확한 증거 의도를 가진 질문으로 변환한다. 이러한 질문은 텍스트 지지, 시각 확인, 교차 모달 설명이라는 세 가지 요구를 포괄한다. 텍스트 지지 질문은 상품 문구나 규격이 리뷰 의견을 지지하는지 확인하고, 시각 확인 질문은 상품 이미지가 관련 구조나 외관을 보여 주는지 확인하며, 교차 모달 설명 질문은 특정 피드백이 실제 제품 문제, 페이지 제시 부족, 사용자 기대 차이 중 어디에 가까운지 판단한다.

검색과 생성 단계에서 시스템은 두 단계 검색을 사용한다 \citep{qdrant2024pdfRetrievalAtScale,qdrant2024multiVectorCourse}. 첫 단계는 벡터 표현을 사용해 상품 텍스트와 상품 이미지에서 후보 증거를 거칠게 회수하고, 두 번째 단계는 재순위 모델을 사용해 후보 집합을 정밀하게 정렬한다. 텍스트 경로는 주로 상품명, 카테고리, 설명, 규격, 명시적 약속을 검증하는 데 사용되고, 이미지 경로는 주로 구조, 외관, 색상, 국소 세부 사항을 검증하는 데 사용된다. 최종 보고서는 측면 집계를 기반으로 장점, 단점, 쟁점, 적합한 사용 장면, 증거 공백, 개선 제안을 생성하고, claim 수준 귀속을 통해 claim과 리뷰, 상품 텍스트, 상품 이미지 사이의 지지 관계를 표시한다.

시스템 구현은 구조화 객체로 다단계 프로세스를 제약한다 \citep{liu2024structuredOutputBenchmark,geng2023grammarConstrainedDecoding,yao2023react}. 모델 호출은 호환형 언어 모델 게이트웨이를 통해 수행되며, 검색 계층은 로컬 인덱스와 벡터 데이터베이스 두 종류의 백엔드를 지원한다. 공식 실험과 개발 검증을 구분하기 위해 시스템은 실행 정책 제어를 제공한다. 전체 모델 체인이 사용 가능할 때는 실제 벡터 인코더, 이미지 인코더, 재순위 모델을 사용하고, 외부 능력이 사용 불가능하며 정책이 허용할 때는 휴리스틱 경로로 후퇴하여 프로세스가 실행 가능하도록 한다. 이 설계는 모델 능력 부족과 프로세스 설계 실패를 혼동하지 않도록 한다.

본 논문의 현재 검증 목표는 고정된 대규모 benchmark의 최종 성능을 보고하는 것이 아니라, 시스템이 증거 기반 VOC 분석을 안정적으로 완료할 수 있는지를 확인하는 것이다. 결과는 현재 구현이 상품 증거 표준화, 리뷰 측면 모델링, 질문 계획, 멀티모달 후보 검색, 후보 재순위, 측면 집계, 보고서 생성, 증거 귀속을 순차적으로 완료할 수 있음을 보여 준다. 하나의 완전한 분석은 최종 자연어 보고서뿐 아니라 리뷰 측면, 질문 의도, 검색 질문, 회수된 증거, 검색 품질, 실행 설정, 측면 집계, 과정 추적, 보고서 귀속 등의 객체를 포함하므로, 사람이 검토하고 오류를 위치시키는 데 도움을 준다.

질문 수준 관련 증거 주석이 있을 때 시스템은 Recall@1, Recall@3, Recall@5, MRR, NDCG@3, NDCG@5 \citep{manning2008introductionIr}를 계산할 수 있다. 구조화 보고서에 대해서는 claim이 리뷰, 상품 텍스트, 상품 이미지에 의해 지지되는지 통계화하며, 여기에는 claim support rate, claim grounded rate, citation precision, citation contradiction rate, 서로 다른 모달리티 경로의 기여 비율이 포함된다. 현재 구현은 166개의 자동화 테스트를 포함하고 모두 통과했으며, 이는 주요 인터페이스, 구조화 중간 표현, 워크플로 단언이 현재 구현에서 일관성을 유지함을 보여 준다.

\subsection{그림 1}

\begin{figure}[H]
\centering
\includegraphics[width=0.92\textwidth]{/home/severin/Codelib/Decathlon\_VOC\_Analyzer/论文\_md形式/图片/图1.png}
\caption{Decathlon VOC Analyzer의 전체 프레임워크. 시스템은 상품 증거 패키지와 리뷰 측면을 중간 표현으로 사용하고, 질문 계획을 통해 주관적 VOC 신호와 멀티모달 상품 증거를 연결하며, 보고서 단계에서 claim 수준 귀속을 수행한다.}
\end{figure}

\subsection{그림 2}

\begin{figure}[H]
\centering
\includegraphics[width=0.92\textwidth]{/home/severin/Codelib/Decathlon\_VOC\_Analyzer/论文\_md形式/图片/图2.png}
\caption{리뷰 측면에서 증거 질의로 이어지는 계획 메커니즘. 측면 객체는 검색에 직접 들어가지 않고, 먼저 서로 다른 증거 의도를 가진 질문으로 분해된 뒤 텍스트 경로, 이미지 경로 또는 이미지-텍스트 결합 경로를 각각 트리거한다.}
\end{figure}

% =========================
% III. 결론
% =========================
\section{III. 결론}
본 논문은 상품 이미지-텍스트 증거와 사용자 리뷰를 정렬하기 위한 증거 기반 멀티모달 VOC 분석 시스템인 Decathlon VOC Analyzer를 제안했다. 시스템은 원시 상품 데이터를 추적 가능한 증거 패키지로 표준화하고, 사용자 리뷰를 측면 수준 VOC 객체로 추출한 뒤, 질문 계획, 멀티모달 후보 검색, 재순위, 보고서 생성, claim 수준 귀속을 통해 구조화된 상품 분석 결과를 생성한다.

본 논문의 핵심 결론은 상품 VOC 분석이 단일 단계 리뷰 요약으로 축소되어서는 안 된다는 것이다. 대신 ``리뷰 측면 모델링 - 증거 질문 계획 - 상품 이미지-텍스트 검색 - 증거 제약 생성''의 다단계 프로세스로 조직되어야 한다. 현재 구현은 실행 가능한 연구 프로토타입을 형성했으며, 구조화 중간 표현, 증거 귀속, replay 메커니즘, 평가 인터페이스에서 일관성을 유지한다. 테스트 결과 현재 구현은 166개의 자동화 테스트를 통과했으며, 이는 시스템 주 체인과 핵심 모듈이 엔지니어링 수준의 안정성을 갖추고 있음을 보여 준다.

향후 연구는 여러 카테고리에 걸친 수작업 주석 benchmark를 보완하고, 다양한 검색 및 재순위 전략을 체계적으로 비교하며, 더 세밀한 시각 grounding을 도입하고, 사람 피드백과 replay 메커니즘을 더욱 긴밀히 연결하는 데 초점을 둘 것이다.

% =========================
% References
% =========================
\bibliographystyle{unsrtnat}
\bibliography{ref}

\end{document}
